\onecolumngrid
\section*{End Matter}
\twocolumngrid

\paragraph*{Hypotheses.}
The Theorem holds under the following, all of which are needed and none of which
is implicit.

\emph{(H1) Locality and no derivatives.} Invariants are Lorentz-invariant
polynomials in the components of $F$ alone. Terms containing $\partial F$ are
outside the class.

\emph{(H2) Self-duality.} $F$ is constrained to the $\dimSelfDualModule$-dimensional
self-dual subspace of $\Lambda^5$, with $\star F = +F$ in signature
$(-,+,\dots,+)$ and $\eps_{01\ldots9}=+1$.

\emph{(H3) Fixed grading.} Degree means polynomial degree in $F$. Degree ten is
the graded piece, not a cumulative count.

\emph{(H4) The reachable sector.} $\Dten$ is the span of degree-ten invariants
generated by the stress-tensor flow of Ref.~\cite{Hutomo:2025chiral} acting on the
free seed theory, closed under the products the flow generates. Its precise
generator list is in the Supplemental Material. A different flow defines a
different $\Dten$ and the theorem must be re-read accordingly.

\emph{(H5) No field redefinitions.} Invariants are compared as polynomials, not
modulo field redefinitions of $F$. Allowing redefinitions could only shrink the
quotient and is a separate question.

\emph{(H6) Arithmetic.} Dimensions are exact ranks over $\Fp$ at the primes
listed in the Supplemental Material. Because the coordinate basis is integral,
each is an unconditional lower bound on the characteristic-zero rank; the
matching upper bounds come from explicit spanning sets, also exact.

\paragraph*{The three representatives.}
Explicit compact representatives of $\Qten$, in the pure-$N$ presentation that
avoids an index ambiguity present in the published expressions, are listed in
Table~I of the Supplemental Material together with their coordinates in the
degree-ten basis and their images in the quotient. Removal minimality is checked
directly: deleting any one representative drops $\dim\Qten$ from \dimQten\ to
\dimQtenMinusOne.

\paragraph*{What is computer-assisted, and how to check it.}
The dimensions in Eq.~\eqref{eq:main} are computer-assisted. They are not opaque:
each is the rank of an explicitly recorded integer matrix, and the repository
ships the matrices, the pivot rows and columns, and a script that recomputes each
rank from the stored data without re-running any search. A reader who distrusts
the enumeration can still check the ranks.

The enumeration itself --- that the recorded candidates exhaust degree ten in the
declared class --- is the step that rests on the search. It is corroborated by an
independent formula-free reverse recovery that reaches the same quotient span,
and by the spinor-side computation, but it is not a symbolic proof and we do not
present it as one.

\paragraph*{What is not claimed.}
The reverse search is bounded and is not an exhaustive enumeration of all
contraction sectors. No generating set for the full invariant ring is presented.
Degree twelve is not classified. The three representatives are a certified choice,
not canonical objects. Agreement at several primes is not a characteristic-zero
identity, and no statement here relies on it being one.
